% Options for packages loaded elsewhere
\PassOptionsToPackage{unicode}{hyperref}
\PassOptionsToPackage{hyphens}{url}
%
\documentclass[
]{article}
\usepackage{amsmath,amssymb}
\usepackage{lmodern}
\usepackage{iftex}
\ifPDFTeX
  \usepackage[T1]{fontenc}
  \usepackage[utf8]{inputenc}
  \usepackage{textcomp} % provide euro and other symbols
\else % if luatex or xetex
  \usepackage{unicode-math}
  \defaultfontfeatures{Scale=MatchLowercase}
  \defaultfontfeatures[\rmfamily]{Ligatures=TeX,Scale=1}
\fi
% Use upquote if available, for straight quotes in verbatim environments
\IfFileExists{upquote.sty}{\usepackage{upquote}}{}
\IfFileExists{microtype.sty}{% use microtype if available
  \usepackage[]{microtype}
  \UseMicrotypeSet[protrusion]{basicmath} % disable protrusion for tt fonts
}{}
\makeatletter
\@ifundefined{KOMAClassName}{% if non-KOMA class
  \IfFileExists{parskip.sty}{%
    \usepackage{parskip}
  }{% else
    \setlength{\parindent}{0pt}
    \setlength{\parskip}{6pt plus 2pt minus 1pt}}
}{% if KOMA class
  \KOMAoptions{parskip=half}}
\makeatother
\usepackage{xcolor}
\setlength{\emergencystretch}{3em} % prevent overfull lines
\providecommand{\tightlist}{%
  \setlength{\itemsep}{0pt}\setlength{\parskip}{0pt}}
\setcounter{secnumdepth}{-\maxdimen} % remove section numbering
\ifLuaTeX
  \usepackage{selnolig}  % disable illegal ligatures
\fi
\IfFileExists{bookmark.sty}{\usepackage{bookmark}}{\usepackage{hyperref}}
\IfFileExists{xurl.sty}{\usepackage{xurl}}{} % add URL line breaks if available
\urlstyle{same} % disable monospaced font for URLs
\hypersetup{
  hidelinks,
  pdfcreator={LaTeX via pandoc}}

\author{}
\date{}

\begin{document}

MAPE, or Mean Absolute Percentage Error, is defined as the average of
the percentage error between the actual and forecasted values. It is
calculated as:

\[
\text{MAPE} = \frac{1}{n} \sum_{t=1}^{n} \left| \frac{A_t - F_t}{A_t} \right| \times 100
\]

There are several advantages to using MAPE:

\begin{enumerate}
\def\labelenumi{\arabic{enumi})}
\item
  It is easy to interpret since it is expressed as a percentage.
\item
  Furthermore, it allows for comparison across different time periods
  and scales
\end{enumerate}

However, the context of predicting the number of passengers that an
airline will carry in the future poses some limitations to the
usefulness of MAPE:

\begin{enumerate}
\def\labelenumi{\arabic{enumi})}
\item
  \textbf{Asymmetric Costs:} MAPE treats over-predictions and
  under-predictions equally, which could be disastrous in this setting.
  Whereas overprediction may lead to some unnecessary operational costs,
  underprediction results in an inadequate fleet capacity, lost revenue
  and disgruntled customers. Therefore, our performance metric needs to
  take this into account.
\item
  \textbf{Inability to Capture Peak Demand:} For fleet management and
  hiring, understanding peak demand is crucial. MAPE averages out the
  errors across all data points, which can mask important peaks or
  troughs in demand.
\item
  \textbf{Instability with Low Values:} MAPE can become
  disproportionately large when actual values approach zero. Given that
  passenger numbers can vary significantly (from 150,000 to over
  6,000,000), low passenger counts in certain months could lead to
  extremely high percentage errors, misleading decision-making.
\end{enumerate}

Given that the fleet requirements are constrained by the total number of
passengers expected while human resources requirements are usually
constrained by peak demand, here are some alternative metrics to
evaluate forecast performance:

\begin{enumerate}
\def\labelenumi{\arabic{enumi})}
\item
  \textbf{Mean Absolute Error (MAE):}

  \begin{itemize}
  \item
    MAE calculates the average of the absolute differences between the
    actual passenger numbers and the forecasted numbers.
  \end{itemize}
\end{enumerate}

\[
\text{MAE} = \frac{1}{n} \sum_{t=1}^{n} |A_t - F_t|
\]

\begin{itemize}
\item
  \textbf{Advantage:} It focuses on passenger count rather than
  percentages, thereby providing a direct assessment of the forecast's
  performance and making it easy for operational teams to understand and
  act upon.
\item
  \textbf{Disadvantage:} It still does not solve the problem of
  asymmetric consequences, since it reports the absolute value of the
  error. Also, it captures only the mean error and not the peak error.
\end{itemize}

\begin{enumerate}
\def\labelenumi{\arabic{enumi})}
\setcounter{enumi}{1}
\item
  \textbf{Weighted MAE:}

  \begin{itemize}
  \item
    We can introduce asymmetry in the calculation of MAE by taking a
    weighted mean of errors, assigning a relatively higher weight for
    underpredictions and a relatively lower weight for overpredictions.
  \item
    However, this loses the interpretability advantage that MAE offered,
    since we have now polluted the MAE with arbitrary weights.
  \end{itemize}
\item
  \textbf{Maximum Absolute Error (MaxAE):}

  \begin{itemize}
  \item
    MaxAE identifies the largest absolute error among all forecasted
    values compared to the actual values.
  \end{itemize}
\end{enumerate}

\[
\text{MaxAE} = \max(|A_t - F_t|) \quad \text{for } t = 1, 2, \ldots, n
\]

\begin{itemize}
\item
  \textbf{Advantage:} It captures the worst-case scenario in terms of
  forecasting error. This is particularly critical for human resource
  planning.
\end{itemize}

\begin{enumerate}
\def\labelenumi{\arabic{enumi})}
\setcounter{enumi}{3}
\item
  \textbf{A} \textbf{Combination of MAE and MaxAE}:

  \begin{itemize}
  \item One can use MAE to plan fleet capacity based on total expected passengers and use MaxAE to ensure operational readiness for peak periods based on worst-case scenarios.
  \end{itemize}
  
\end{enumerate}

\end{document}
